% No preamble. Prefix zif-. The optional fixed-coupling parity refinement stays in research.
\section{Quadratic forcing from a fixed outer field domain}
\label{zif-section}

Fix the composite geometry $R,J$, carrier exponent $a_0=1/16$,
and the jointly chosen, input-independent classical system of
the preceding local-limit theorem. The gauge Lie algebra is
semisimple, in particular it may be $\mathfrak{su}(2)$.
Constants may depend on the fixed geometry and color conventions,
but are uniform in preceding homogeneous depth $M$ and compatible
volume. The input of the quantum step in this section is zero.

\subsection{One outer construction and its restrictions}
\label{zif-outer-section}

The classical precision, stationary-map, pivot and base-interaction
coefficients have independent outer product domains. After $R,J$
are fixed, choose one radius $\rho>0$ on which all their strict
support and Gaussian-neighborhood budgets hold with a fixed
margin. No inherited activity is evaluated in this construction.
With the actual retained link/color window of threshold $T$, write
\begin{align}
 \mathcal G_{M,g,T}(0)
   &=\Delta_M+\mathcal Q_{{\rm det},M}
                         +\mathcal L_{M,g,T},\notag\\
 \mathcal L_{M,g,T}(w)
   &=\log\mathbb E_{C_M(w)}
             [\chi_T(\xi)e^{V_{\rm base}(w,\xi;g)}].
 \label{zif-zero-decomposition}
\end{align}
The covariance, stationary branch and charts are independent of
$g$. Field-independent powers of $g$, flat determinants, Haar
constants and the compact gauge-insertion normalization remain
in the separately tracked vacuum factor.

The actual outer-domain construction gives an anchored family
with
\begin{align}
 \|\mathcal G_{M,g,T}(0)\|_{a_0,\rho}
 &\le H(\rho,g,T),\notag\\
 H(\rho,g,T)
 &:=(C_\Delta+C_{\rm det})\rho
       +2S_1\bigl[e^3\tau(T)+e^{v(g,T)}v(g,T)\bigr],\notag\\
 v(g,T)&\le C_{\rm base}(gT^3+gT+g^2T^2),
 \label{zif-outer-bound}
\end{align}
with harmless fixed enlargements of the constants if needed.
Choose $T=T(g)\ge1$ and $g_0>0$ so that every strict base and
bad-block budget holds on $\rho$ for $0<g\le g_0$ and
\begin{equation}
 \sup_{M,\,0<g\le g_0}H(\rho,g,T(g))\le H_*<\infty.
 \label{zif-uniform-outer}
\end{equation}
A sufficiently large fixed multiple of
$\sqrt{\log(1/g)}$, with a fixed lower threshold, is one
permitted choice after reducing $g_0$. This is a bound on a
coefficient family, not on an extensive total.

For $0<b<\rho$, fix these outer data and restrict their
coefficients. The zero-input integral on $b$ is exactly the
restriction of \eqref{zif-zero-decomposition}: its chart,
covariance, window $T(g)$, measure factors and logarithm branch
are unchanged. A field radius enters the estimates, not the
defining integral. The infinite statement uses the same
outer-domain local-limit system followed by restriction.
Unrelated subsequences at different radii are not identified.
Two alternative valid carrier assignments for that same finite
integral also give the same normalized class by the established
local-limit theorem.

\subsection{Remove the total constant and linear sectors}
\label{zif-force-section}

\begin{theorem}[Uniform quadratic zero-input force]
\label{zif-theorem}
Under \eqref{zif-uniform-outer}, for every $b\le\rho/2$ the
zero-input correction has a representative of its
vacuum-normalized class satisfying
\begin{equation}
 \|[\mathcal G_{M,g,T(g)}(0)]\|_{\mathcal Q,a_0,b}
       \le C_F b^2,
 \qquad C_F=\frac{2H_*}{\rho^2}.
 \label{zif-force-bound}
\end{equation}
The constant is independent of $b,M$ and volume. The same
representative estimate holds before taking a local limit if
the hypotheses are imposed on the actual finite regulators.
\end{theorem}

\begin{proof}
The actual zero-input total is globally adjoint invariant:
the real retained-window change of variables, equivalently
the local gauge theorem of the common-domain integral, gives
this identity. Its derivative at zero in each link slot is
an invariant linear functional on the adjoint Lie algebra.
It vanishes on commutators, and
$[\mathfrak g,\mathfrak g]=\mathfrak g$; hence its total
linear jet is zero. This reasoning passes from finite
regulators to compact-field derivatives on the chosen limit.
Individual carrier coefficients need not have zero linear
jets or be separately gauge invariant.

Let $P_n$ be homogeneous phase projection in all background
coordinates of one carrier. Cauchy's formula on radii below
$\rho$, followed by their increasing limit, gives
\begin{equation}
 \|P_n E_H\|_b\le(b/\rho)^n\|E_H\|_\rho.
 \label{zif-phase-bound}
\end{equation}
Replace the outer family by
\begin{equation}
 E_H^{\ge2}:=\sum_{n\ge2}P_nE_H.
 \label{zif-representative}
\end{equation}
It retains every declared carrier and represents the same
normalized total: constants are normalized away and the
sum of the linear sectors is zero. Sum
\eqref{zif-phase-bound} first over degree for each carrier,
then over anchored carriers. This yields
\begin{equation}
 \|E^{\ge2}\|_{a_0,b}
 \le\frac{(b/\rho)^2}{1-b/\rho}\|E\|_{a_0,\rho}
 \le\frac{2H_*}{\rho^2}b^2.
 \label{zif-geometric-bound}
\end{equation}
There is no exchange of a root supremum with an order sum.
The argument is identical at finite volume and on the
normally convergent infinite family.
\end{proof}

For a Lie algebra with an abelian factor, the adjoint argument
alone does not remove a linear total; its vanishing would be
an additional hypothesis. Vacuum removal here changes only
the normalized class. It does not discard the corresponding
physical scalar normalization.

Let the actual full reference satisfy
$\beta(S_{MR})=1$ and
$\|[S_{MR}]\|_{\mathcal Q,b}\le C_Sb^2$, as obtained
from its fixed outer classical representation. Then
\begin{equation}
 P_{MR}=I-S_{MR}\beta,
 \qquad \|P_{MR}\|\le p:=1+C_\beta C_S
 \label{zif-projection}
\end{equation}
has a bound independent of $b$. Therefore
\begin{align}
 \|P_{MR}\mathcal G_{M,g,T(g)}(0)\|_{\mathcal Q,b}
      &\le C_0b^2,\qquad C_0=pC_F,\notag\\
 |\beta(\mathcal G_{M,g,T(g)}(0))|
      &\le C_\beta C_F.
 \label{zif-extracted-beta}
\end{align}
The $b^{-2}$ factor in the norm of $\beta$ is present;
the proven $b^2$ forcing cancels it. The finite representative
estimate does not by itself define a finite extraction
functional: a finite $\beta_N$ must separately have the
corresponding bounded normalization on the actual reference.

\subsection{Order of choices and the limits of the conclusion}
\label{zif-scope-section}

The outer radius and bound are chosen once after $R,J$.
The final inner radius may subsequently be reduced to meet
any fixed finite block of classical transport constraints.
For a prescribed family ceiling $\bar\epsilon$, the inherited
quantum difference coefficient is proportional to $gT(g)/b$.
A sufficient budget, with the constants of the common-domain
theorem, is
\begin{equation}
 \frac{16C_TC_RgT(g)}b\,\bar\epsilon
       \le\frac{\eta_*}{128e}.
 \label{zif-inherited-budget}
\end{equation}
The denominator $128e$ is the stronger two-input comparison
budget used for the quotient Lipschitz estimate. It is met by
reducing $g_0$ after the positive $b$ and
$\bar\epsilon$ have been fixed. Thus a ceiling of the form
$\bar\epsilon\le\eta_*/(\text{constant}\,a_Q)$,
where $a_Q=16C_TC_RgT/b$, creates no radius cycle.
The base and bad-block budgets remain separate requirements.

If all intermediate extracted amplitudes have norm at most
$C_{\rm int}b^2$, the bounded classical transport, the force
bound and the fixed-coupling quantum Lipschitz bound give
an unprojected one-step amplitude bounded by a constant
times $b^2$. Its Maxwell increment is consequently bounded
independently of $b$. This provides neither its sign nor
control of infinitely many accumulated increments.

Finally, $\Delta_M$ and $\mathcal Q_{{\rm det},M}$ have
no small factor $g$. Their removal of constant and linear
sectors yields $O(b^2)$, not $O(g^2)$. A bound of the latter
kind for the full force would require a separate proved
cancellation or subtraction. Nothing here is a contraction
in the joint coupling--activity variables, a large-field
estimate, or an identification of the local sector with the
full Wilson measure.
